\chapter{Теория}
\label{ch:intro}

\section*{Дискретные сигналы}

Дискретный сигнал - последовательность дискретных отсчетов сигнала

$$X(n) = {1,2,3,4}, n = {-1, 0, 1, 2}$$

Здесь $X(n)$ - значения сигнала в моменты времени $n$. \\

Единичный импульс - базовый импульс для описания произвольного дискретного сигнала. Аналог дельта-функции в дискретной области.
Применяется для описания произвольных дискретных сигналов. \\

Произвольная последовательность отсчетов может быть описана как сумма взвешанных и задержанных единичных импульсов.

$$X(n) = \sum_{k=-\infty}^{\infty}X(k)\delta(n-k)$$

Здесь $X(k)$ отвечает за взвешанность единичного импульса (масштабирует его). $\delta(n-k)$ - сам единичный импульс, отвечает за
задержанность.

\section*{Дискретные системы}

Дискретные системы выполняют преобразование дискретных сигналов

\begin{figure}[H]
    \centering
    \includegraphics[width=1.0\textwidth]{dis_sys.png}
    \caption{Дискретная система}
\end{figure} 

Здесь $F{*}$ - оператор преобразования (какой-то алгоритм).\\

$$h(n) = F{\delta(n)}$$

Здесь $h(n)$ (импульсная характеристика) - реакция системы на единичный импульс. \\

\subsection{FIR(КИХ) системы}

FIR(КИХ) системы - дискретные системы с конечной импульсной характеристикой. В такой системе ИХ состоит из конечного кол-ва отсчетов.

$$y(n) = \sum_{k=M}^{N}b_kX(n-k)$$

$$h(n) = \sum_{k=M}^{N}b_n\delta(n-k)$$

Для систем КИХ выходное значение в любой момент времени n зависит только от текущего и предыдущего x(n)

\subsection{IIR(БИХ) системы}

Данные системы описываются разностными уравнением:

$$\sum_{k=0}^{N}a_ky(n-k) = \sum_{k=0}^{N}b_kx(n-k)$$

Если $a_0$ = 1

$$y(n) = \sum_{k=0}^{N}b_kX(n-k) - \sum_{k=1}^{N}a_kX(n-k)$$

В такой системе выход соединен обратной связью со входом и текущий выходной сигнал зависит от предыдущих выходных сигналов. Благодаря
обратной связи создается эффект бесконечной ИХ.

\section{Дискретная свертка}

Произвольная последовательность отсчетов может быть описана как сумма взвешанных и задержанных единичных импульсов. Значит, на выходе
системы получим:

$$y(n) = F{\sum_{k=-\infty}^{\infty}X(k)\delta(n-k)} = \sum_{k=-\infty}^{\infty}F{X(k)\delta(n-k)} = $$

$$\sum_{k=-\infty}^{\infty}X(k)F{\delta(n-k)} = \boxed{\sum_{k=-\infty}^{\infty}X(k)h(n-k)}$$

В рамке находится дискретная свертка.


\section{Способы вычисления дискретной свертки}

Представлять вычисление дискретной свертки можно разными способами, но все они дают один результат

\subsection{Прямое суммирование}

\begin{figure}[H]
    \centering
    \includegraphics[width=1.0\textwidth]{first.png}
    \caption{Вариант вычисления дискретной свертки}
\end{figure} 

Каждый элемент отсчет S(n), поступающий на вход ФНЧ с ИХ h(n), будет иметь иметь вес S(n) и задержку h(n), т.е каждый S(n) будет
повторяться 3 раза, т.к h(n) имеет длину 3. Дублируем так каждый отсчет, а потом суммируем все полученные отсчеты в каждый момент времени. 

\subsection{Зеркальное отображение ИХ}

\begin{figure}[H]
    \centering
    \includegraphics[width=1.0\textwidth]{second.png}
    \caption{Вариант вычисления дискретной свертки}
\end{figure} 

Разворачиваем ИХ в обратную сторону и двигаем навстречу к семплам и перемножаем. Далее полученные значения нужно просуммировать.

\endinput